\subsection{La personnalisation de l'apprentissage}
\label{subsec:personnalisation_apprentissage}

Les capacités de génération contextuelle décrites en~\ref{subsec:agents_scriptes_generatifs} ouvrent une forme de personnalisation inaccessible aux systèmes à règles fixes. Un LLM adapte en temps réel la formulation, le registre et le niveau de détail au fil de l'échange, sans que le concepteur ait à anticiper chaque cas. Cette adaptation sémantique produit un effet asymétrique : la personnalisation du contenu par un LLM --- phrases et récits ajustés au profil de l'apprenant --- améliore la motivation intrinsèque, mais n'affecte pas la performance d'apprentissage mesurée par des tests de compréhension \citep{leong2024}. L'adaptation agit donc sur l'intérêt perçu, pas sur le traitement cognitif du contenu. Ce pattern --- perception améliorée, apprentissage inchangé --- se reproduit dans les travaux présentés ci-après.

Au-delà du contenu, la personnalisation peut porter sur l'agent lui-même. L'utilisation de personnalités admirées comme instructeurs vidéo améliore la motivation et les émotions positives, sans modifier les scores d'apprentissage \citep{pataranutaporn2022}. Le passage au format interactif accentue cet effet : lorsque les apprenants dialoguent avec l'instructeur virtuel via un chat textuel, l'engagement augmente par rapport au visionnage passif, les scores d'apprentissage restant équivalents \citep{prasongpongchai2024}. L'interactivité ajoute ainsi une dimension que la vidéo seule ne capture pas. L'effet atteint son maximum lorsque la connexion repose sur une relation déjà établie. La comparaison entre un chatbot textuel, une mascotte 3D et un deepfake de l'enseignant réel des élèves montre que la familiarité préexistante avec la personne représentée augmente la confiance et la compétence perçue de l'agent \citep{tan2025}. La gradation est nette : image d'une figure admirée, puis interaction avec cette figure, puis relation préexistante avec la personne imitée --- chaque étape renforce la perception sans que l'apprentissage mesuré ne suive. Les comparaisons directes entre agents génératifs et instructeurs humains confirment ce constat : l'équivalence rapportée porte sur des mesures immédiates dans des contextes contrôlés, non sur une reproduction du processus d'enseignement humain \citep{leiker2023, lim2024}.

Le mécanisme sous-jacent relève d'un transfert de crédibilité : l'apprenant reporte sur l'agent la confiance qu'il accorde déjà à la personne représentée, comme le montre la gradation décrite ci-dessus. L'apparence humaine agit en outre comme un prédicteur de l'utilité perçue, un effet médié par l'interactivité attendue plutôt que par la ressemblance elle-même \citep{lin2024}. Cette dissociation entre perception et apprentissage, déjà documentée pour les agents classiques (cf.~\ref{subsec:incarnation_agence_sociale}), pose un problème de conception. L'attrait visuel de l'agent peut détourner des ressources attentionnelles du contenu vers l'interface, un phénomène observé y compris avec des instructeurs humains filmés (cf.~\ref{subsec:incarnation_agence_sociale}). Cette tension entre engagement perçu et traitement cognitif s'accentue lorsque l'agent incarne une figure historique.

Des institutions culturelles déploient déjà des agents conversationnels incarnant des figures disparues. Au Dalí Museum de St. Petersburg, l'installation \textit{Ask Dalí}\footnote{\url{https://thedali.org/exhibit/ask-dali/}} permet aux visiteurs de converser avec une reconstruction de Salvador Dalí via un téléphone-homard, objet iconique de son \oe{}uvre surréaliste (voir Figure~\ref{fig:ask_dali}). Le système, alimenté par GPT-4 et une synthèse vocale entraînée sur des archives audio, produit des réponses reflétant la personnalité de l'artiste. Au Musée d'Orsay, l'expérience \textit{Hello Vincent}\footnote{\url{https://www.musee-orsay.fr/fr/articles/numerique-bonjour-vincent-275618}} propose une interaction avec Van Gogh appuyée sur un corpus de près de 900 lettres validé par un historien de l'art (voir Figure~\ref{fig:vangogh_museum}). Ces déploiements marquent le passage du prototype de recherche au grand public. Ils révèlent aussi une tension : le cadre muséal vise la médiation culturelle et l'engagement émotionnel, tandis que le cadre scolaire exige l'exactitude factuelle et l'apprentissage conceptuel.

\begin{figure}[htbp]
    \centering
    \includegraphics[width=0.7\textwidth]{images/ch2/dali.png}
    \caption{Installation \textit{Ask Dalí} au Dalí Museum de St. Petersburg. Les visiteurs conversent avec une reconstruction IA de Salvador Dalí via un téléphone-homard. La citation murale --- « I do not understand why, when I ask for a grilled lobster in a restaurant, I am never served a cooked telephone » --- illustre l'humour surréaliste que le système cherche à reproduire.}
    \label{fig:ask_dali}
\end{figure}

\begin{figure}[htbp]
    \centering
    \includegraphics[width=0.7\textwidth]{images/ch2/vangogh.png}
    \caption{Représentation numérique de Vincent van Gogh dans l'expérience \textit{Hello Vincent} au Musée d'Orsay. Le visiteur converse avec l'artiste pendant qu'il peint son célèbre \textit{Champ de blé aux corbeaux}, l'IA s'appuyant sur un corpus de près de 900 lettres.}
    \label{fig:vangogh_museum}
\end{figure}

Plusieurs travaux explorent cette personnification dans un cadre expérimental. Des participants adultes conversant avec des reconstructions de Leonardo Da Vinci ou Murasaki Shikibu rapportent une motivation et une efficacité d'apprentissage perçue supérieures à la lecture seule \citep{pataranutaporn2023} (voir Figure~\ref{fig:living_memories}). Des collégiens dialoguant avec Napoléon ou Alexandre le Grand montrent une amélioration de l'empathie historique et des capacités d'auto-réflexion \citep{kim2025}. Les variables mesurées --- motivation, empathie, engagement --- relèvent toutefois du registre affectif. L'apprentissage factuel, la compréhension conceptuelle et la rétention à long terme restent largement inexplorés, conformément au pattern identifié plus haut. L'effet optimal survient lorsque l'interaction complète les sources textuelles plutôt que de s'y substituer \citep{pataranutaporn2023}. La valeur pédagogique de ces agents semble résider dans la complémentarité, non dans la substitution.

\begin{figure}[htbp]
    \centering
    \includegraphics[width=0.8\textwidth]{images/ch2/living-memories-talking.png}
    \caption{Interface du projet \textit{Living Memories} \citep{pataranutaporn2023}. L'utilisateur converse avec une reconstruction de Leonardo Da Vinci, qui se présente et répond aux questions. Le système combine GPT-3 et un extracteur sémantique pour générer des réponses cohérentes avec le personnage historique.}
    \label{fig:living_memories}
\end{figure}

La personnalisation par l'apparence soulève également des questions que la littérature commence à explorer \citep{danry2022}. Le consentement de la personne représentée ne peut être obtenu lorsqu'il s'agit de figures historiques décédées. Les effets d'une exposition prolongée à des représentations synthétiques sur la perception de l'authenticité restent inconnus. La confusion entre l'agent et la personne qu'il imite constitue un risque, en particulier pour de jeunes apprenants. Au-delà de ces enjeux liés à l'apparence de l'agent, la fiabilité du contenu qu'il produit constitue un problème distinct, examiné dans la section suivante.

